% DELETE -- os/samir/cmd/mutate.c m_delete (default the current record;
% FOR without a scope = ALL; marks with 0x2A, written at once).
% Gate: test-interp-replace, test-dbf-mutate.
\begin{XbaseEntry}{DELETE}{command}
Marks records of the database in USE for deletion.
\Syntax
\SyntaxLine{\kw{DELETE} \opt{\xp{scope}} \opt{\kw{FOR} \xp{condition}}
  \opt{\kw{WHILE} \xp{condition}}}
\Defaults
Without a scope or a condition, \cmd{DELETE} marks the current record only.
With a \kw{FOR} condition and no scope, the scope is \kw{ALL}.
\Usage
A record marked for deletion stays in the database.  \cmd{LIST} and
\cmd{DISPLAY} show it with an asterisk in front of its record number, and the
function \fn{DELETED()} is true for it.  \cmd{RECALL} removes the mark;
\cmd{PACK} removes the marked records from the database for good.

While SET DELETED is ON, the commands that process records leave out the
records marked for deletion.

The mark is written to the disk at once.
\Examples
\DotExample{examples/c-delete}
\SeeAlso{PACK, RECALL, SET DELETED, ZAP, DELETED()}
\end{XbaseEntry}

